\oldpage[III]{190}
\subsection{Caso de los anillos locales noetherianos}
\label{subsection:III.7.5}

\begin{env}[7.5.1]
\label{III.7.5.1}
Sean $A$ un anillo local noetheriano y $\mathfrak m$ su ideal
maximal. Para todo $A$-módulo $M$, designemos por $\widehat M$ su
completado separado para la topología $\mathfrak m$-ádica, el cual es
isomorfo a
\[
  \varprojlim_n
  \bigl(M\otimes_A(A/\mathfrak m^{n+1})\bigr)
  =
  \varprojlim_n
  \bigl(M/\mathfrak m^{n+1}M\bigr).
\]
Sea $T$ un funtor covariante aditivo de $\cat{Ab}_A$ en
$\cat{Ab}$. Los homomorfismos canónicos
\sref{III.7.2.3.1}
\[
  T(M)\otimes_A(A/\mathfrak m^{n+1})
  \longrightarrow
  T\bigl(M\otimes_A(A/\mathfrak m^{n+1})\bigr)
\]
forman evidentemente un sistema proyectivo de $A$-homomorfismos. Por
paso al límite dan, por tanto, un $\widehat A$-homomorfismo funtorial
en $M$
\[
  \widehat{T(M)}
  \longrightarrow
  \varprojlim_n T(M_n),
  \tag{7.5.1.1}\label{III.7.5.1.1.eq}
\]
donde se ha puesto
\[
  M_n=M\otimes_A(A/\mathfrak m^{n+1}),
  \qquad
  A_n=A/\mathfrak m^{n+1}.
\]
\end{env}

\begin{proposition}[7.5.2]
\label{III.7.5.2}
Sean $A$ un anillo local noetheriano de ideal maximal $\mathfrak m$,
$k=A/\mathfrak m$ su cuerpo residual y $T$ un funtor covariante
aditivo de $\cat{Ab}_A$ en $\cat{Ab}$, semiexacto y que conmuta con
los límites inductivos. Supongamos además que, para todo $A$-módulo
de tipo finito $M$, $T(M)$ es un $A$-módulo de tipo finito y que el
homomorfismo canónico \eqref{III.7.5.1.1.eq} es un isomorfismo. En
estas condiciones, las propiedades siguientes son equivalentes:
\begin{enumerate}
  \item[(a)] $T$ es exacto por la derecha;
  \item[(b)] para todo $n$, el funtor $N\mapsto T(N)$ es exacto por
    la derecha en la categoría de los $A_n$-módulos de tipo finito;
    equivalentemente, $T$ es exacto por la derecha en la categoría
    de los $A$-módulos de longitud finita;
  \item[(c)] el homomorfismo canónico
    \[
      T(\varepsilon):T(A)\longrightarrow T(k)
    \]
    es sobreyectivo;
  \item[(d)] para todo $n$ suficientemente grande, el homomorfismo
    canónico
    \[
      T(A_n)\longrightarrow T(k)
    \]
    es sobreyectivo.
\end{enumerate}
\end{proposition}

\begin{proof}
Es claro que (a) implica (b). Demostremos el recíproco. Sea
$u:M\to N$ un epimorfismo de $A$-módulos; debemos probar que $T(u)$
es un epimorfismo. Como $T$ conmuta con los límites inductivos y los
límites inductivos filtrantes son exactos en la categoría de los
módulos, basta considerar el caso en que $M$ y $N$ son de tipo
finito.

Los módulos $T(M)$ y $T(N)$ son entonces de tipo finito. Como $A$ es
noetheriano y local, basta probar que
\[
  \widehat{T(u)}:
  \widehat{T(M)}
  \longrightarrow
  \widehat{T(N)}
\]
es sobreyectivo \sref[0]{0.7.3.5} y \sref[0]{0.6.4.1}. Por hipótesis,
$\widehat{T(M)}$ y $\widehat{T(N)}$ son, respectivamente,
\[
  \varprojlim_n T(M_n)
  \quad\hbox{y}\quad
  \varprojlim_n T(N_n).
\]
Así, $\widehat{T(u)}$ es el límite inverso del sistema proyectivo de
homomorfismos
\[
  T(u\otimes 1_{A_n}):T(M_n)\longrightarrow T(N_n).
\]
La condición (b) significa precisamente que estos homomorfismos son
sobreyectivos. Además, $T(M_n)$ es un $A_n$-módulo de tipo finito y
$A_n$ es artiniano. Se sigue que $\widehat{T(u)}$ es sobreyectivo
\sref[0]{0.13.1.2} y \sref[0]{0.13.2.2}.

Es claro que (a) implica (c), y (c) implica (d). Finalmente,
\sref{III.7.2.7} muestra que (b) y (d) son equivalentes, pues $T$ es
semiexacto en $\cat{Ab}_{A_n}$. Esto termina la demostración.
\end{proof}

\begin{corollary}[7.5.3]
\label{III.7.5.3}
Bajo las hipótesis generales de \sref{III.7.5.2}, si $T(k)=0$,
entonces $T(M)=0$ para todo $A$-módulo $M$.
\end{corollary}

\begin{proof}
\oldpage[III]{191}
Como $k$ es el único $A$-módulo simple, de
\eqref{III.7.3.5.4.eq} se deduce que $T(E)=0$ para todo $A$-módulo
de longitud finita $E$. Si $M$ es de tipo finito, entonces
\[
  \widehat{T(M)}
  \simeq
  \varprojlim_n T(M_n).
\]
Cada $M_n$ es de longitud finita, de modo que
$\widehat{T(M)}=0$. Por hipótesis, $T(M)$ es de tipo finito y se
inyecta en $\widehat{T(M)}$ \sref[0]{0.7.3.5}; por consiguiente,
$T(M)=0$. Finalmente, para un $A$-módulo arbitrario $M$, el módulo
$T(M)$ es el límite inductivo de los $T(N_\alpha)$, donde los
$N_\alpha$ recorren los submódulos de tipo finito de $M$. Esto
termina la demostración.
\end{proof}

\begin{proposition}[7.5.4]
\label{III.7.5.4}
Sean $A$ un anillo local noetheriano de ideal maximal $\mathfrak m$,
$k=A/\mathfrak m$ su cuerpo residual y $T_\bullet$ un funtor
homológico de $\cat{Ab}_A$ en $\cat{Ab}$ que conmuta con los límites
inductivos. Supongamos además que, para todo $i$ y todo $A$-módulo
de tipo finito $M$, $T_i(M)$ es de tipo finito y el homomorfismo
canónico
\[
  \widehat{T_i(M)}
  \longrightarrow
  \varprojlim_n T_i(M_n)
\]
es biyectivo. Para un entero dado $p$, las condiciones siguientes
son equivalentes:
\begin{enumerate}
  \item[(a)] $T_p$ es exacto;
  \item[(b)] $T_p$ es exacto por la derecha y $T_p(A)$ es un
    $A$-módulo libre;
  \item[(c)] los homomorfismos canónicos
    \[
      T_{p+1}(A)\longrightarrow T_{p+1}(k),
      \qquad
      T_p(A)\longrightarrow T_p(k)
    \]
    son sobreyectivos;
  \item[(d)] para todo $n$, los homomorfismos canónicos
    \[
      T_{p+1}(A_n)\longrightarrow T_{p+1}(k),
      \qquad
      T_p(A_n)\longrightarrow T_p(k)
    \]
    son sobreyectivos;
  \item[(e)] para todo $n$, el funtor $N\mapsto T_p(N)$ es exacto
    en la categoría de los $A_n$-módulos de tipo finito.
\end{enumerate}
\end{proposition}

\begin{proof}
Por \sref{III.7.3.3}, la condición (a) equivale a decir que
$T_{p+1}$ y $T_p$ son exactos por la derecha. Como $T_\bullet$ es
un funtor homológico en $\cat{Ab}_{A_n}$, el mismo razonamiento que
en \sref{III.7.3.1} muestra que (e) equivale a decir que $T_p$ y
$T_{p+1}$ son exactos por la derecha en la categoría de los
$A_n$-módulos de tipo finito. Se deduce de \sref{III.7.5.2} que (a)
y (e) son equivalentes; la misma proposición da también la
equivalencia de (a), (c) y (d).

Finalmente, todo módulo plano de tipo finito sobre el anillo local
$A$ es libre (Bourbaki, \emph{Algèbre commutative}, cap.~II,
\S~3, n\textsuperscript{o}~2, cor.~2 de la prop.~5). La equivalencia
de (a) y (b) resulta entonces de \sref{III.7.3.1} y
\sref{III.7.3.3}.
\end{proof}

\begin{corollary}[7.5.5]
\label{III.7.5.5}
Supongamos verificadas las hipótesis generales de
\sref{III.7.5.4}.
\begin{enumerate}
  \item[(i)] Si $T_p(k)=0$, entonces $T_p=0$, el funtor $T_{p+1}$
    es exacto por la derecha y $T_{p-1}$ es exacto por la izquierda.
  \item[(ii)] Si $T_{p-1}(k)=T_{p+1}(k)=0$, entonces $T_p$ es
    exacto, el homomorfismo canónico
    \[
      T_p(A)\otimes_A M
      \longrightarrow
      T_p(M)
    \]
    es biyectivo y $T_p(A)$ es un $A$-módulo libre.
\end{enumerate}
\end{corollary}

\begin{proof}
La aserción (i) se deduce inmediatamente de \sref{III.7.5.3}, pues
$T_p$ es semiexacto; la última aserción resulta de la definición de
funtor homológico. Teniendo en cuenta \sref{III.7.3.3}, (i) da de
inmediato las dos primeras aserciones de (ii), y el hecho de que
$T_p(A)$ sea libre resulta de \sref{III.7.5.4}.
\end{proof}

\begin{corollary}[7.5.6]
\label{III.7.5.6}
Supongamos verificadas las hipótesis generales de
\sref{III.7.5.4}, y supongamos además que $A$ es un anillo de
valuación discreta.
\begin{enumerate}
  \item[(i)] El funtor $T_p$ es exacto por la derecha si y solo si
    $T_{p-1}(A)$ es un $A$-módulo libre.
  \item[(ii)] El funtor $T_p$ es exacto si y solo si $T_p(A)$ y
    $T_{p-1}(A)$ son $A$-módulos libres.
\end{enumerate}
\end{corollary}

\oldpage[III]{192}
\begin{proof}
Es claro que (ii) implica (i) (cf.\ \sref{III.7.3.3}). Para probar
(i), sea $f$ un generador del ideal maximal de $A$, es decir, un
uniformizador de $A$. Un $A$-módulo de tipo finito $M$ es libre,
equivalentemente plano, si y solo si la homotecia
\[
  h_f:M\longrightarrow M,\qquad x\longmapsto fx,
\]
es inyectiva, pues en la situación presente esto equivale a decir
que $M$ es sin torsión \sref[0]{0.6.3.4}. La sucesión exacta
\[
  0
  \longrightarrow A
  \xrightarrow{h_f}
  A
  \longrightarrow k
  \longrightarrow 0
\]
da la sucesión exacta de homología
\[
  T_p(A)
  \longrightarrow T_p(k)
  \longrightarrow T_{p-1}(A)
  \xrightarrow{h_f}
  T_{p-1}(A).
\]
Así, $T_{p-1}(A)$ es libre si y solo si
$T_p(A)\to T_p(k)$ es sobreyectivo, y la conclusión resulta de
\sref{III.7.5.2}.
\end{proof}

\begin{remark}[7.5.7]
\label{III.7.5.7}
En virtud de \sref{III.7.4.7}, las hipótesis generales de
\sref{III.7.4.4} implican que el funtor homológico $T_\bullet$
definido por \eqref{III.7.4.1.1.eq} satisface las hipótesis generales
de \sref{III.7.5.4}. En este caso, \sref{III.7.5.6} está, por tanto,
contenido en \sref{III.7.4.3}.
\end{remark}
